\documentclass[11pt,a4paper]{article}

% -------------------------------------------------
% HSLU-like Proposal Template (DSPRO1 FS26)
% -------------------------------------------------

% ---------- Encoding & Typography ----------
\usepackage[T1]{fontenc}
\usepackage[utf8]{inputenc}
\usepackage[english]{babel}
\usepackage{lmodern}
\usepackage{microtype}

% ---------- Page Layout ----------
\usepackage[a4paper,margin=25mm]{geometry}
\usepackage{setspace}
\setstretch{1.15}
\setlength{\parindent}{0pt}
\setlength{\parskip}{0.5em}

% ---------- Colors ----------
\usepackage{xcolor}
\definecolor{hsluRed}{HTML}{E30613}
\definecolor{hsluDark}{HTML}{222222}
\definecolor{hsluGray}{HTML}{666666}
\definecolor{lineGray}{HTML}{D9D9D9}
\definecolor{boxGray}{HTML}{F5F5F5}

% ---------- Graphics & Tables ----------
\usepackage{graphicx}
\usepackage{tabularx}
\usepackage{array}
\usepackage{booktabs}
\usepackage{longtable}

% ---------- Headers/Footers ----------
\usepackage{fancyhdr}
\pagestyle{fancy}
\fancyhf{}
\lhead{\small Hochschule Luzern \textbar\ Informatik}
\rhead{\small DSPRO1 FS26 Proposal}
\cfoot{\small \thepage}
\renewcommand{\headrulewidth}{0.4pt}
\renewcommand{\headrule}{\hbox to\headwidth{\color{lineGray}\leaders\hrule height \headrulewidth\hfill}}
\renewcommand{\footrulewidth}{0pt}

% ---------- Section Style ----------
\usepackage{titlesec}
\titleformat{\section}
  {\Large\bfseries\color{hsluDark}}
  {\thesection}{0.6em}{}
\titleformat{\subsection}
  {\normalsize\bfseries\color{hsluDark}}
  {\thesubsection}{0.6em}{}
\titlespacing*{\section}{0pt}{1.2em}{0.4em}
\titlespacing*{\subsection}{0pt}{0.8em}{0.2em}

% ---------- Hyperlinks ----------
\usepackage[hidelinks]{hyperref}

% ---------- Utility ----------
\newcommand{\blankline}{\vspace{0.2em}{\color{lineGray}\hrule}\vspace{0.55em}}
\newcommand{\fieldlabel}[1]{\textbf{#1}\vspace{0.15em}}
\newcommand{\hsluwordmark}{%
  {\fontsize{22}{24}\selectfont\bfseries\textcolor{black}{HSLU}}%
}
\newcommand{\proposaltype}{%
  {\fontsize{13}{15}\selectfont\bfseries Project Proposal}
}
\newcommand{\moduleline}{%
  {\fontsize{11}{13}\selectfont DSPRO1}
}

\begin{document}

% =========================
% Cover / Title Area
% =========================
\begin{minipage}[t]{0.62\textwidth}
    \hsluwordmark\\[-1mm]
    {\small\color{hsluGray}Hochschule Luzern}\\[-0.5mm]
    {\small\color{hsluGray}Departement Informatik}
\end{minipage}
\hfill
\begin{minipage}[t]{0.35\textwidth}
    \raggedleft
    \proposaltype\\[1mm]
    \moduleline
\end{minipage}

\vspace{8mm}
{\color{hsluGray}\rule{\textwidth}{1.2pt}}

\vspace{8mm}

\begin{center}
    {\LARGE\bfseries\color{hsluDark}Proposal for Data Science Project}\\[1mm]
    {\large\color{hsluGray}Semester: FS26}
\end{center}

\vspace{8mm}

% =========================
% Metadata Box
% =========================
\colorbox{boxGray}{%
\parbox{\dimexpr\linewidth-2\fboxsep\relax}{
\vspace{1mm}
\renewcommand{\arraystretch}{1.35}
\begin{tabularx}{\linewidth}{>{\bfseries}p{0.26\linewidth} X}
Project Title & Predicting Apartment Rental Prices in Switzerland Using Machine Learning \\
Course / Module & DSPRO1 FS26 \\
Team Members & Elias Martinelli, Timo Schlumpf \\
Supervisor / Coach & Elena Nazarenko \\
Date of Submission & 03.03.2026 \\
\end{tabularx}
\vspace{1mm}
}
}

\vspace{1em}

% =========================
% Main Sections
% =========================
\section{Short Project Description}

The Swiss housing market is highly location-dependent and rental prices can vary substantially across cities, cantons, and neighborhoods. For prospective tenants, it is often difficult to assess whether the rent of a listed apartment is justified compared to similar properties currently available on the market.

This project aims to develop a supervised machine learning model to predict monthly apartment rental prices in Switzerland based on publicly available real estate listing data. The objective is to support data-driven rent estimation by providing a transparent and reproducible price prediction approach. The machine learning task is a regression problem, where the target variable is the monthly apartment rent in CHF.

Multiple regression models will be evaluated, including a simple baseline model, Linear Regression, and Random Forest Regressor. This allows us to compare predictive performance between a simple interpretable model and a more flexible tree-based model. Our hypothesis is that location-related variables together with structural apartment features, such as living area and number of rooms, contain enough signal to predict rental prices with practically useful accuracy, and that Random Forest will outperform Linear Regression.

The primary stakeholders are prospective tenants who want to evaluate whether a rental offer is fair relative to the market. A secondary stakeholder group could be landlords or property owners who want a rough market-based estimate of an appropriate rental price. The expected impact is an improved understanding of the main drivers of apartment rents and a practical model that helps assess whether a listed apartment is fairly priced.


\section{Data Description}

The primary dataset will consist of publicly available apartment rental listings from the Swiss real estate market collected via web scraping from one or more major Swiss real estate platforms. We expect to collect several thousand rental listings, depending on data availability and platform structure. If the available Swiss dataset is too limited for robust model training and evaluation, an additional housing dataset may be used during development for methodological testing. However, the main analytical focus remains on the Swiss rental market.

The target variable is the monthly apartment rent in CHF, treated as a continuous variable in a regression setting. The dataset is expected to contain structured, tabular listing information. Numerical variables may include living area in square meters, number of rooms, floor, and year built. Categorical variables may include canton, city, ZIP code, property type, and apartment characteristics such as balcony, parking, or furnished status.

Initial preprocessing will include handling missing values, removing duplicate listings, cleaning inconsistent categorical labels, and encoding categorical variables. Depending on the final feature set, scaling of numerical variables may be applied for linear models. Additional checks will be performed to identify implausible values and outliers, such as unrealistic apartment sizes or extreme rental prices. All predictor variables will be selected such that they are available at listing time in order to avoid data leakage.

Model selection will be performed using a hold-out split and cross-validation. This will allow a structured comparison of models while keeping the final evaluation separate from model tuning.


\section{Cloud Service Integration}

Model development, preprocessing, and training will be conducted locally in Python using standard machine learning libraries such as pandas, scikit-learn, and matplotlib. GitHub will serve as the central platform for version control, collaboration, and documentation throughout the project.

No dedicated cloud deployment or external cloud computing service is currently planned, since the main focus of the project is on data collection, preprocessing, model development, and evaluation rather than on production deployment. This setup keeps the technical scope aligned with the course objectives while ensuring transparency and reproducibility.

References and literature sources will be managed using Zotero.


\section{Kanban Tool}

The project will be managed using GitHub Projects integrated with the project repository. The Kanban board will include the columns Backlog, Ready, In Progress, Review, and Done.

Tasks will be created as GitHub issues and moved across columns as development progresses through the project. Planned tasks include data scraping, data cleaning, exploratory data analysis, feature engineering, model training, evaluation, and documentation. Each task will be assigned to a responsible team member. This setup ensures structured task tracking, clear task ownership, and transparent project progress.


\section{Experiment Tracking Tool}

MLflow will be used for experiment tracking and model management. The tool will log model parameters, selected features, preprocessing configurations, dataset version identifiers, evaluation metrics, and relevant artifacts such as trained models and result plots. This ensures reproducibility and systematic comparison of model variants.

The main evaluation metrics will be Mean Absolute Error (MAE), Root Mean Squared Error (RMSE), and the coefficient of determination ($R^2$). MAE provides an intuitive measure of average prediction error in CHF, RMSE penalizes larger prediction errors more strongly, and $R^2$ indicates how much of the variance in rental prices can be explained by the model.

As a baseline, a simple model predicting the mean rental price of the training data will be implemented first. Its performance will then be compared with Linear Regression and Random Forest Regressor in order to assess whether more flexible machine learning models provide a meaningful improvement over simpler approaches.

\end{document}
